\section*{Chapitre \ref{ch:b3:molecular-evolution} --- Évolution moléculaire et phylogénomique}

\begin{solution}{exo:b3:molecular-evolution:1}
$k = \mu$~: le \omterm{thm:b3:molecular-evolution:neutral}{taux de substitution} neutre par site et par génération
égale le taux de mutation. À chaque génération $2N\mu$ nouvelles
mutations neutres apparaissent et chacune a une probabilité $1/2N$ de se
fixer~; une population plus grande fait plus de mutations et donne à
chacune une chance proportionnellement plus petite, et les deux effets
s'annulent exactement.
\end{solution}

\begin{solution}{exo:b3:molecular-evolution:2}
\omterm{def:b3:genomics:comparative}{Orthologues}~: des gènes de deux espèces descendus d'un seul gène de leur
ancêtre commun --- la $\beta$-globine humaine et la $\beta$-globine de
souris. \omterm{def:b3:genomics:comparative}{Paralogues}~: des gènes d'un même génome descendus d'une
duplication --- l'$\alpha$- et la $\beta$-globine humaines, ou $\beta$ et
$\gamma$. \omterm{def:b3:molecular-evolution:duplication}{Pseudogène}~: une copie délabrée qui ne code plus de protéine
--- le gène $\psi\beta$ du groupe de la $\beta$-globine humaine, codons
stop compris.
\end{solution}

\begin{solution}{exo:b3:molecular-evolution:3}
$\omega$ compare le rythme des substitutions qui changent un acide aminé
à celui des substitutions silencieuses, qui est la référence neutre.
$0.02$~: \omterm{def:b3:molecular-evolution:dnds}{sélection purificatrice} forte, \qty{98}{\%} des changements
d'acide aminé retirés --- l'actine, l'\omterm{def:b3:chromatin-epigenetics:nucleosome}{histone} H3. $1.0$~: aucune
contrainte --- un \omterm{def:b3:molecular-evolution:duplication}{pseudogène}. $2.5$~: \omterm{def:b3:molecular-evolution:dnds}{sélection positive}, changements
d'acide aminé favorisés --- le sillon de liaison au peptide du CMH,
l'enveloppe du VIH, le lysozyme des ruminants.
\end{solution}

\begin{solution}{exo:b3:molecular-evolution:4}
Les protéines comparées entre mammifères changent à raison d'environ une
substitution par site et par milliard d'années~; sur un génome d'un
milliard de sites codants et une génération de quelques années, cela
fait de l'ordre d'une substitution fixée par génération, soit une tous
les deux ans. Le coût de la sélection de Haldane autorise environ une
substitution sélectionnée toutes les trois cents générations. Les
rythmes diffèrent d'un facteur cent ou plus, de sorte que presque tous
les changements fixés doivent être neutres.
\end{solution}

\begin{solution}{exo:b3:molecular-evolution:5}
Neutre~: $1/2N = 10^{-4}$. $s = 0.001$, $4Ns = 20$~: $(1 - \mathrm{e}^{-0.002})/
(1 - \mathrm{e}^{-20}) = 0.002$. $s = 0.01$~: $0.0198$. $s = -0.001$~: $(1 -
\mathrm{e}^{0.002})/(1 - \mathrm{e}^{20}) = 0.002\,\mathrm{e}^{-20} =
4\times 10^{-12}$. Aucune n'est effectivement neutre~: il faudrait pour
cela $|4Ns|
\lesssim 1$, c'est-à-dire $|s| \lesssim 5\times 10^{-5}$~; un coefficient
de sélection d'un dixième de pour cent est déjà décisif dans une
population de cinq mille.
\end{solution}

\begin{solution}{exo:b3:molecular-evolution:6}
Non corrigé~: $T = 0.12/(2\times 2\times 10^{-9}) = \qty{30}{Myr}$.
Jukes--Cantor~: $d = -\tfrac{3}{4}\ln(1 - 0.16) = 0.131$, $T =
\qty{33}{Myr}$. La correction atteint un facteur deux quand $-\tfrac{3}{4}
\ln(1 - 4p/3) = 2p$, vers $p \approx 0.6$~: à ce stade la différence
brute a perdu l'essentiel de son information.
\end{solution}

\begin{solution}{exo:b3:molecular-evolution:7}
$p_{N} = 7/700 = 0.010$, $p_{S} = 10/200 = 0.050$, $\omega = 0.2$~:
\omterm{def:b3:molecular-evolution:dnds}{sélection purificatrice} qui retire quatre cinquièmes des changements
d'acide aminé. Avec $35$ et $10$~: $p_{N} = 0.05$, $\omega = 1$~: aucune
contrainte décelable --- un \omterm{def:b3:molecular-evolution:duplication}{pseudogène}, ou un gène où une \omterm{def:b3:molecular-evolution:dnds}{sélection
positive} à certains sites compense une \omterm{def:b3:molecular-evolution:dnds}{sélection purificatrice} à
d'autres.
\end{solution}

\begin{solution}{exo:b3:molecular-evolution:8}
$k = d/2T = 0.012/(1.3\times 10^{7}) = 9.2\times 10^{-10}$ par site et
par an~; $\mu = 25k = 2.3\times 10^{-8}$ par site et par génération~;
$6\times
10^{9}\times 2.3\times 10^{-8} \approx 140$ mutations nouvelles par
génome diploïde et par génération. Le séquençage direct de parents et
d'enfants en trouve environ soixante-dix~: l'estimation par la
divergence est gonflée par le polymorphisme déjà présent dans la
population ancestrale, qui ajoute quelques centaines de milliers de
générations à la séparation apparente.
\end{solution}

\begin{solution}{exo:b3:molecular-evolution:9}
$\mathrm{e}^{-T/2N} = \mathrm{e}^{-1} = 0.37$~; fraction discordante $\tfrac{2}{3}
\times 0.37 = 0.25$. Pour \qty{10}{\%}~: $\mathrm{e}^{-T/2N} = 0.15$, $T =
2N\ln(1/0.15) = \num{190000}$ générations.
\end{solution}

\begin{solution}{exo:b3:molecular-evolution:10}
Quand $s \to 0$~: $1 - \mathrm{e}^{-2s} \approx 2s$ et $1 - \mathrm{e}^{-4Ns}
\approx 4Ns$, rapport $1/2N$. Pour $4Ns \gg 1$ le dénominateur vaut $1$ et
le numérateur $\approx 2s$. Pour $s < 0$ avec $4N|s| \gg 1$~: numérateur $1 -
\mathrm{e}^{2|s|} \approx -2|s|$, dénominateur $1 - \mathrm{e}^{4N|s|} \approx
-\mathrm{e}^{4N|s|}$, donc $P \approx 2|s|\,\mathrm{e}^{-4N|s|}$. Cent fois
sous la neutre dans $N = 10^{4}$~: $2|s|\,\mathrm{e}^{-4\times 10^{4}|s|} =
5\times 10^{-7}$~; $|s| = 1.6\times 10^{-4}$ donne $3.2\times 10^{-4}\times
\mathrm{e}^{-6.4} = 5.3\times 10^{-7}$. Donc $|s| \approx 1.6\times 10^{-4}$,
$4N|s| \approx 6.5$~: un désavantage d'un centième de pour cent suffit,
chez dix mille individus, à rendre la fixation cent fois plus rare que
le hasard.
\end{solution}

\begin{solution}{exo:b3:molecular-evolution:11}
Avec le point de séparation à la distance $a$ d'$O$ sur les deux
chemins, $a + m =
0.30$, $a + h = 0.24$, $m + h = 0.20$~: $m - h = 0.06$, donc $m = 0.13$, $h
= 0.07$, rapport $1.9$. Une horloge stricte par génération donnerait le
rapport des nombres de générations, $25/0.5 = 50$. Le $1.9$ observé
signifie que la souris n'a accumulé que le double du changement humain
par an malgré cinquante fois plus de générations~: le taux de mutation
par génération doit être quelque $25$ fois plus élevé chez l'humain ---
plus de divisions cellulaires germinales par génération --- de sorte que
par an les deux lignées ne diffèrent que d'un facteur deux.
\end{solution}

\begin{solution}{exo:b3:molecular-evolution:12}
Les substitutions créant un stop arrivent à $1000\times 2\times 10^{-9}\times
0.04 = 8\times 10^{-8}$ par an~: la première est attendue au bout
d'environ \qty{12}{Myr} (les décalages du cadre divisent cela à peu près
par deux). Dans cette fenêtre une mutation bénéfique, une sur $10^{3}$ à
$10^{5}$ de toutes les mutations, a bien moins de chances d'arriver
qu'un codon stop, et même quand elle arrive elle ne se fixe qu'avec la
probabilité $2s$. La copie est donc d'ordinaire morte avant que la
sélection ait pu lui trouver un usage --- la demi-vie observée des
duplicats chez les vertébrés est de quelques millions d'années.
\end{solution}

\begin{solution}{pb:b3:molecular-evolution:1}
\textbf{1.} $k = 0.012/(2\times 6.5\times 10^{6}) = 9.2\times 10^{-10}$ par
site et par an~; $\mu = 25k = 2.3\times 10^{-8}$ par site et par génération.
\textbf{2.} $6\times 10^{9}\times 2.3\times 10^{-8} \approx 140$ mutations
nouvelles par génome diploïde~; \qty{1.5}{\%} d'entre elles, environ
$2$, dans du codant.
\textbf{3.} Par génome haploïde, $3\times 10^{9}\times 2.3\times 10^{-8}
\approx 70$ substitutions fixées par génération, contre le $1/300$ de
Haldane~: vingt mille fois plus que ce que la sélection pourrait mener,
de sorte que l'écrasante majorité est neutre.
\textbf{4.} $2N\mu = 10^{5}\times 2.3\times 10^{-8} = 2.3\times 10^{-3}$
mutations nouvelles par site et par génération dans la population~;
$1/2N =
10^{-5}$ d'entre elles se fixeront.
\textbf{5.} $4N = \num{200000}$ générations, \qty{5}{Myr} --- presque
aussi longtemps que le temps écoulé depuis la séparation.
\textbf{6.} La divergence compte les mutations apparues le long des deux
lignées séparées depuis leur ancêtre commun~; chacune se fixe dans sa
propre lignée, et le rythme à long terme vaut $\mu$ quel que soit le
temps que chacune met, puisque les mutations apparaissent régulièrement
et que le retard ne fait que différer le flux sans le changer. (La seule
correction concerne le polymorphisme présent à la séparation, qui rend
l'ancêtre commun de deux allèles plus vieux que la séparation des
espèces.)
\textbf{7.} Neutre $10^{-5}$~; $s = 0.001$~: $0.002$~; $s = 0.01$~: $0.0198$~;
$s = 0.1$~: $1 - \mathrm{e}^{-0.2} = 0.18$.
\textbf{8.} $s = 1/4N = 5\times 10^{-6}$~: en dessous la dérive noie la
sélection et la mutation se comporte comme neutre~; au-dessus la
sélection gouverne les chances.
\textbf{9.} $s = -0.001$~: $P = 0.002\,\mathrm{e}^{-200}$, nulle à toutes
fins utiles --- $10^{-85}$ de la neutre. $s = -10^{-5}$, $4N|s| = 2$~: $P =
2\times 10^{-5}/(\mathrm{e}^{2} - 1) = 3.1\times 10^{-6}$, \qty{31}{\%} de
la neutre~: les mutations faiblement délétères se fixent bel et bien.
\textbf{10.} Substitutions bénéfiques par site et par génération~: $2N\times
10^{-5}\mu\times 2s = 10^{5}\times 10^{-5}\times 0.02\,\mu = 0.02\mu$~;
neutres~: $\approx\mu$. Fraction adaptative $0.02/1.02 \approx \qty{2}{\%}$.
Bien que chaque mutation bénéfique ait deux mille fois plus de chances
de se fixer qu'une neutre, elles sont si rares qu'une substitution sur
cinquante seulement est adaptative~: la sélection agit et se voit à peine.
\textbf{11.} $p_{N} = 46/920 = 0.050$, $p_{S} = 84/280 = 0.30$. Corrigés~:
$d_{N} = -\tfrac{3}{4}\ln(1 - 0.0667) = 0.052$, $d_{S} = -\tfrac{3}{4}
\ln(1 - 0.40) = 0.38$. $\omega = 0.14$.
\textbf{12.} \omterm{def:b3:molecular-evolution:dnds}{Sélection purificatrice} qui retire environ \qty{86}{\%} des
changements d'acide aminé --- un gène typique. $k = d_{S}/2T = 0.38/(1.8\times
10^{8}) = 2.1\times 10^{-9}$ par site et par an~: le double de la valeur
humain--chimpanzé, comme on l'attend d'une comparaison qui inclut la
lignée des rongeurs à l'horloge rapide~; le même ordre de grandeur.
\textbf{13.} $d = -\ln(1 - 0.55) = 0.80$ par site~; $T = 0.80/(2\times
10^{-9}) = \qty{400}{Myr}$.
\textbf{14.} Deux chaînes différentes font le tétramère
$\alpha_{2}\beta_{2}$, dont la liaison coopérative de l'oxygène, l'\omterm{prop:b3:structural-biology:mechanism}{effet
Bohr} et le contrôle allostérique par le 2,3-bisphosphoglycérate dépendent
de l'interface entre sous-unités dissemblables~; et le groupe $\beta$ a
pu ensuite se diversifier en chaînes embryonnaires, fœtales et adultes
d'affinités différentes.
\textbf{15.} $2kT = 2\times 8\times 10^{-9}\times 4\times 10^{7} = 0.64$
substitution par site~; avec la saturation entre vingt acides aminés la
différence brute vaut environ
$0.95(1 - \mathrm{e}^{-0.64/0.95}) \approx 0.47$. À \qty{500}{Myr},
$d = 8$~: chaque site a changé bien des fois et la différence brute se
tient à son plafond de \qty{95}{\%}, sans porter aucune information sur
le temps.
\textbf{16.} $10^{-11}\times 10^{9}\times 100\times 2 = 2$ substitutions
entre deux lignées sur 100 sites en un milliard d'années. Pour une
séparation de \qty{10}{Myr} l'espérance vaut $0.02$~: presque toujours
aucune différence, donc aucune résolution.
\textbf{17.} $1200\times 2\times 10^{-9}\times 0.04 = 9.6\times 10^{-8}$
par an~: le premier stop est attendu au bout d'environ \qty{10}{Myr}.
\textbf{18.} Doubler le génome entier double tous les gènes à la fois, de
sorte que la stœchiométrie des protéines qui interagissent est
préservée et qu'aucun déséquilibre de dosage ne sélectionne contre les
copies~; un duplicat en tandem isolé change la dose d'un seul gène et est
souvent délétère, ou aussitôt redondant et libre de se délabrer.
\textbf{19.} $\mathrm{e}^{-\num{80000}/\num{100000}} = \mathrm{e}^{-0.8} = 0.45$.
\textbf{20.} $\tfrac{2}{3}\times 0.45 = 0.30$~: cohérent avec les
\qty{30}{\%} observés.
\textbf{21.} La moitié chacune, \qty{15}{\%} de tous les arbres groupant
l'humain avec le gorille et \qty{15}{\%} le chimpanzé avec le gorille.
Un net excès de l'un indiquerait un flux de gènes entre ces deux espèces
après leur séparation, que la dérive ne peut produire.
\textbf{22.} $N = \num{10000}$~: $\mathrm{e}^{-4} = 0.018$, discordance
\qty{1.2}{\%}. Les \qty{30}{\%} observés exigent donc une population
ancestrale de quelque cinquante mille --- plusieurs fois l'effectif
efficace des humains d'aujourd'hui.
\textbf{23.} La concaténation suppose que chaque gène a l'arbre des
espèces~; avec le \omterm{prop:b3:molecular-evolution:ils}{tri incomplet des lignées} les gènes ont de nombreux
arbres, et pour des branches internes courtes l'arbre de gène le plus
fréquent peut différer de l'arbre des espèces, de sorte que plus de
données rendent la mauvaise réponse plus assurée. Une méthode fondée sur
le coalescent calcule, pour chaque \omterm{prop:b3:molecular-evolution:ils}{arbre d'espèces} candidat, la
distribution attendue des \omterm{prop:b3:molecular-evolution:ils}{arbres de gènes}, et retient l'\omterm{prop:b3:molecular-evolution:ils}{arbre d'espèces}
qui explique le mieux les fréquences observées.
\textbf{24.} Le \omterm{prop:b3:molecular-evolution:ils}{tri incomplet des lignées} s'est produit dans la
population ancestrale commune de tous les humains modernes, de sorte
qu'il rendrait toutes les populations humaines également apparentées aux
Néandertaliens. Un excès chez les non-Africains signifie que des allèles
néandertaliens sont entrés après le départ d'Afrique des ancêtres des
non-Africains~: un métissage, il y a quelque cinquante mille ans.
\textbf{25.} $\mu \approx 2.3\times 10^{-8}$ par site et par génération~;
$\omega \approx 0.14$~; environ \qty{30}{\%} des \omterm{prop:b3:molecular-evolution:ils}{arbres de gènes} sont
discordants avec l'arbre des espèces.
\end{solution}
